% s-022-p022-m002; original PDF 22, 23
\Aphorism{s-022-p022-m002}{13}
Es giebt noch andere Arten von Natur als jene supergescheidte Mechanik, deren Minimumseigenschaften das Entzücken des achtzehnten Jahrhunderts erregten, eines sokratischen Jahrhunderts, das auch den Menschen als Mechanismus verstand und in Fragen der Moral, Bildung, Erziehung ohne die \Sperr{Instincte} rechnete. Unter diesen Instincten ist aber einer, man nenne ihn den Instinct der Gegensätzlichkeit, Zwecklosigkeit, Kraftverschwendung, der vielleicht die uralte missverständliche Abtrennung des Menschen von der Natur in einem präcisen Sinne erneuern lässt, als eine letzte und grundlegende Anthropodicee: Menschlichkeit als der Ausnahmefall von Natürlichkeit, der \Sperr{nicht} nach Minimumsgesetzen lebt, der \Sperr{nicht} mit kleinstem Kraftaufwand und kleinstem Zwange ein kleinstes Mass von Unlust zu verwirklichen trachtet, sondern dem es erst im klingenden Spiel der grossen Gefahren und Contraste wohl wird. Der Mensch, das \Foreign{animal} überhaupt, strebt nicht nach Glück und entscheidet sich nicht in der Linie des „geringsten Widerstandes“; mit solch einem Nützlichkeits-Calcul im Instinct wird man bourgeois, Thier, Pflanze, Molekül, aber nicht Mensch. Die aufsteigende biologische Reihe, der wir angehören, beweist einen Willen zur Differentiirung, der mit einer Art Heroismus der dumpfen hedonistischen Schwerkraft widerstrebt; ohne diesen „Zug nach oben“ wäre das ursprüngliche Gleichgewicht des Anorganischen entweder nie verlassen oder (falls sich zufällig organische Keime in den tellurischen Brei verirrt hätten) sehr bald wieder erreicht worden. So wenig wie das Organische dürfte man die \Sperr{Cultur} innerhalb jenes philiströsen Mechanismus suchen, der auf Kraftersparniss und geringste Reibung hinarbeitet; jede geistige Epoche redet von ant-eudämonistischen Grundströmungen, von einer gewissen Gleichgültigkeit gegen Lust oder Unlust, von überschäumender Kampf- und Leidensbereitschaft. Jedes Kunstwerk ist ein Ringen mit Elementargeistern, ein Überwältigen des widerspenstigen Stoffs, der gestaltlos bleiben will, ein Heraufklingen des Kosmos aus dem Chaos; jeder Künstler stellt, solange er am Werke ist, einen Herakles vor, der sich für die \Greek{ἀρετή} und nicht für die \Greek{ἡδονή} entschieden hat.
\par

